

\fsection[\textcolor{waveBlue2}{Ejercicio 7}]{\boldit{\textcolor{waveBlue2}{7.}}{
Pregunta a alguna persona mayor de tu entorno sobre su estilo de vida y sus hábitos de consumo cuando era joven y compáralos con los de la juventud actual ( tipos de productos y servicios que consumía , frecuencia de la compra , ropa , ocio , reciclaje , reparación de productos , utilización de bolsas y envases de plástico , etc. ) . ¿ Qué diferencias encuentras entre estos estilos de vida respecto al impacto ambiental que producen ?
}}

\begin{solution}
\begin{enumerate}
    \item \textbf{Consumo de productos y servicios}
    \begin{itemize}
        \item \textbf{Antes:} Compraban menos productos, de uso duradero; predominaba la ropa de fabricación local y electrodomésticos reparables.
        \item \textbf{Ahora:} Consumo frecuente de productos tecnológicos, moda rápida y servicios digitales.
    \end{itemize}

    \item \textbf{Frecuencia de la compra}
    \begin{itemize}
        \item \textbf{Antes:} Compras más espaciadas; se planeaban con antelación.
        \item \textbf{Ahora:} Compras frecuentes, muchas veces impulsivas, incluso online.
    \end{itemize}

    \item \textbf{Ropa y ocio}
    \begin{itemize}
        \item \textbf{Antes:} Ropa duradera, ocio local o al aire libre, actividades comunitarias.
        \item \textbf{Ahora:} Moda rápida, ocio digital, viajes frecuentes y consumo de entretenimiento global.
    \end{itemize}

    \item \textbf{Reciclaje y reparación}
    \begin{itemize}
        \item \textbf{Antes:} Reparaban electrodomésticos y ropa; reciclaje limitado pero uso consciente de recursos.
        \item \textbf{Ahora:} Menos reparación, más desecho; reciclaje presente pero a veces insuficiente.
    \end{itemize}

    \item \textbf{Uso de bolsas y envases de plástico}
    \begin{itemize}
        \item \textbf{Antes:} Bolsas reutilizables, envases retornables o de vidrio.
        \item \textbf{Ahora:} Plástico de un solo uso, envases desechables, aunque con campañas de reducción y reciclaje.
    \end{itemize}

    \item \textbf{Diferencias en impacto ambiental}
    \begin{itemize}
        \item La juventud actual genera mayor huella ambiental por consumo frecuente, moda rápida, tecnología desechable y uso de plásticos.
        \item La generación anterior tenía un estilo de vida más sostenible de facto: productos duraderos, reparables y menor consumo global.
    \end{itemize}
\end{enumerate}
\end{solution}
